\section{Direct calibration of the abstention rule}
\label{sec:calibration}

The design document fixed a detectable halving of per-cell output ($s=0.5$),
coverage of at least $90\%$, and discrimination of at least $80\%$, defined as
the probability that the interval excludes zero at that effect. The proposed
calibration chooses the smallest tested count meeting both criteria for an
unqualified estimate, and the smallest meeting coverage alone for a
wide-interval estimate. The panel, labelling axes and falsification rule are
frozen in configuration and checked against exact values in the test suite.
These safeguards record the specification; they do not establish its adequacy.
Appendix~\ref{app:methods} specifies the cohorts and within-arm cell bootstrap,
which is distinct from patient resampling for population inference.

The substrate comprises ten paired patients in GSE132465 (SMC) and six in
GSE144735 (KUL3). We analyse GUCA2A in the differentiated population under the
fixed lineage labelling. Each synthetic draw uses a two-patient source pool
and 2{,}000 cells per arm, with the reference mature fraction fixed at $0.40$.
The nominal 95\% percentile interval resamples mature cells within each arm
while holding fractions fixed. Its target is conditional intrinsic uncertainty;
it does not propagate uncertainty in the fractions or estimate uncertainty
across a population of patients.

\paragraph{Counts and candidate cutpoints.}
A mature cell belongs to the cell type whose per-cell output defines the
estimand; $n$ counts these cells in the diseased arm. No observed mature cells
means an undefined sample mean, whereas positive but small counts raise an
uncertainty question. An absent underlying population has no conditional-mean
referent. We denote the joint coverage-and-discrimination candidate by $ok$
and the coverage-only candidate by $wide$; the committed rule sets them to
50 and 20. The draw pool specifies eligible source cells, the grid specifies
tested counts, and bins aggregate interval outcomes across counts. We examine
each choice below.

\paragraph{The draw pool changes the result.}
The generator filters eligibility by held-out patient, pooling cells from both
tissues before drawing either synthetic arm. The design document did not
specify tissue eligibility. We therefore evaluate both this implemented pool
and a reference-tissue-only pool. With pooled cells, no count bin meets both
criteria: discrimination reaches at most $0.78$ across 13 seeds at 50
replicates, and $0.75$ across eight seeds at 500. The reference-only pool does
return candidates, but does not recover the committed count of 50
(Figure~\ref{fig:calibration}).

\begin{figure}[htbp]
\centering
\includegraphics[width=\linewidth]{figures/full/fig1_calibration.pdf}
\caption{Direct coverage and discrimination evaluation of the cell-count rule.
Dashed lines mark the pre-specified targets; shading marks the committed
three-state rule. Left: the implemented pooled-cell generator returns no
qualifying cutpoint at any tested count. Right: the reference-only generator
returns a candidate of 90 on seven of eight seeds in the plotted filled-grid
run, rather than the committed 50. Shaded curve bands span the eight seeds;
they are not patient-level confidence intervals. The historical three-grid
comparison below is followed by a controlled union-grid analysis.}
\label{fig:calibration}
\end{figure}

\paragraph{Reported cutpoints depend on binning and grid resolution.}
The committed grid varies mature-cell fractions at 2{,}000 cells per arm,
reaching counts $\{0,20,40,100,200,400,800\}$. It has no point between 40 and
100. On the reference pool it returns $ok=100$ on all eight seeds at 500
replicates, and twelve of thirteen at 50.

Coverage and discrimination are aggregated into twelve log-spaced count bins.
Each reported count is the median of the counts landing in its bin, rather
than a directly estimated crossing. Adding nine fractions to the grid moves
$ok$ from 100 to 90 while retaining the same selected bin, $(76,107]$:
the committed grid contributes $\{100\}$ there, while the filled grid
contributes $\{80,90,100\}$. The additional counts also change the evidence
aggregated within the bin. The newly populated bin $(55,76]$ has median
discrimination $0.789$, below target.

A third grid, reaching below 20, returns $ok=70$ on all eight seeds.
Its fractional count summaries are bin medians, not fractional cells.
The reported crossing brackets and candidates differ:

\begin{center}\small
\begin{tabular}{@{}lccc@{}}
\toprule
 & \textbf{committed} & \textbf{filled} & \textbf{dense} \\
\midrule
crossing bracketed to & $(40,100]$ & $(65,90]$ & $(42.5,70]$ \\
reported $ok$ & $100$ & $90$ & $70$ \\
reported $wide$ & $40$ & $45$ & $42$ \\
replicates per grid point & $500$ & $500$ & $200$ \\
\bottomrule
\end{tabular}
\end{center}

\paragraph{A controlled comparison isolates grid and binning effects.}
These historical brackets describe adjacent reported count summaries, not
confidence intervals. Because their budgets, bin boundaries and streams also
differed, they established configuration sensitivity but did not isolate its
source. We therefore simulated the union of counts once for each cohort, draw
pool and eight seeds (200 outer replicates and 200 inner bootstrap draws), then
formed every grid and binning summary from those same draws. With common fixed
bins, the first cohort returns 100 on the committed grid and 70 on both expanded
grids on all eight seeds. Adaptive bins instead return 65 on seven of eight
filled-grid seeds (90 on one) and 70 on every dense-grid seed. The original
100/90/70 sequence therefore combined missing count support with binning and
new random draws; it was not a pure resolution effect. In the second cohort,
every grid and binning scheme agrees within a seed: six seeds return 400 and two
return 800. Pooled draws still return no qualifying cutpoint in either cohort.
The largest binomial Monte Carlo standard error is 0.035, conditional on the
fixed cohorts; it is not patient-sampling uncertainty.

\begin{table}[htbp]
\centering\small
\caption{Joint coverage-and-discrimination candidates from the controlled
union-grid run. Entries give candidates with seed frequencies in parentheses
where they vary; otherwise all eight seeds agree. Every grid summary reuses the
same draws at shared counts. ``None'' means no tested count or bin qualifies.}
\label{tab:controlledgrid}
\begin{tabular}{@{}llccc@{}}
\toprule
cohort / pool & grid & fixed bins & adaptive bins & per count \\
\midrule
SMC / reference & committed & $100$ & $100$ & $100$ \\
 & filled & $70$ & $65\ (7),\ 90\ (1)$ & $50$--$80$ \\
 & dense & $70$ & $70$ & $45$--$80$ \\
KUL3 / reference & all three & $400\ (6),\ 800\ (2)$ & same & same \\
Both / pooled & all three & None & None & None \\
\bottomrule
\end{tabular}
\end{table}

The per-count summaries in Table~\ref{tab:controlledgrid} supply a useful
check on the bin medians. In SMC, the filled grid's first qualifying actual
count is 50 on two seeds, 60 on five and 80 on one; the dense grid adds a
candidate of 45 on one seed. Thus a stable binned candidate of 70 is not
evidence that 70 is a uniquely determined biological threshold. A bin pools
outcomes over several count settings, whereas a per-count rate addresses only
that setting. Fixed bins hold the boundaries constant, but adding settings can
still change both a bin's median count and its aggregate performance.

These candidates select the first qualifying evaluated point, without fitting
a monotone response or requiring all larger counts to qualify. Their adjacent
crossing brackets describe grid support, not confidence intervals for a
threshold. Conditional Monte Carlo errors quantify rate precision; they do not
account for searching several counts and cannot certify a selected rule's
out-of-sample coverage. The controlled comparison resolves the historical
confounding of grid choice with random draws, while leaving the decision rule
itself unvalidated.

\paragraph{Simulation precision is not patient-level uncertainty.}
Enumerating all 45 and 15 two-patient holdout pairs separates patient omission
from repeated simulation. Keying streams to the sorted pair identifiers keeps
every retained pair fixed. Raising the inner bootstrap budget from 200 to 1{,}000
and 5{,}000 leaves the first cohort's candidate at 70, but moves the second's
400 to 300 and back to 400. Omitting one source patient changes the joint
coverage-and-discrimination candidate for none of ten patients in the first
cohort and three of six in the second (KUL01, KUL30 and KUL31). For the
coverage-only candidate, the corresponding counts are five of ten and zero of
six. These are cohort-conditional influence checks, not intervals for a larger
patient population. We therefore report candidates, not an endorsed calibrated
cutpoint. Appendix~\ref{app:gridcaveats} records denominator and coincident-cutpoint
issues, and Appendix~\ref{app:guards} examines small-sample interval calibration.

\begin{table}[htbp]
\centering\small
\caption{Reference-only sensitivity with every two-patient holdout pair
represented once per count and effect. The three inner budgets reuse identical
outer samples. Omission counts compare with the baseline at 200 inner draws;
all pairs involving the omitted patient are removed and retained pairs are
unchanged. These are sensitivity summaries within the observed cohorts.}
\label{tab:sensitivity}
\begin{tabular}{@{}lcrrrcc@{}}
\toprule
 & pairs & \multicolumn{3}{c}{joint candidate, inner draws} &
 \multicolumn{2}{c}{omissions changing candidate} \\
\cmidrule(lr){3-5}\cmidrule(lr){6-7}
cohort & & $200$ & $1{,}000$ & $5{,}000$ & joint & coverage only \\
\midrule
SMC & $45$ & $70$ & $70$ & $70$ & $0/10$ & $5/10$ \\
KUL3 & $15$ & $400$ & $300$ & $400$ & $3/6$ & $0/6$ \\
\bottomrule
\end{tabular}
\end{table}

The sensitivity windows cover counts $40,50,60,70,80,90,100,110$ in SMC and
$200,300,400,600,800$ in KUL3. These per-count summaries use a fixed pair
schedule, so they are a separate experiment from the eight-seed union grid.
Omitting KUL01 or KUL30 moves the joint candidate from 400 to 300; omitting
KUL31 moves it to the lowest tested count, 200, with its lower boundary
unobserved. The coverage-only baselines likewise occur at the lowest tested
counts, 40 and 200. Their stability cannot locate a lower crossing outside
those windows. Enumerating pairs improves transparency about source-patient
influence, but overlapping pairs do not create additional independent patients.

\paragraph{A second cohort gives the same calibration failure.}
The second Lee cohort contains six patients, compared with ten in the first.
On the committed and filled grids, pooled draws return no qualifying cutpoint
at either replicate count. Discrimination alone reaches $0.80$ at 50
replicates but at most $0.66$ at 500. Reference-only draws yield $ok=400$ or
800, eight to sixteen times the committed 50. At 50 replicates, eleven of
thirteen seeds return 800; at 500, six of eight return 400. This cohort has
only fifteen distinct holdout pairs, and its largest grid point requests 800
mature cells per arm from a pool of 844 distinct mature cells.

Thus neither the recovery curve nor seed stability justified the committed
rule. The algebraic equality replicates across both cohorts, while direct
calibration remains sensitive to pool, grid and cohort. On real samples,
abstention dominates at the finest annotation resolution
(Appendix~\ref{app:modal}); the calibration failure does not license replacing
those unavailable contrasts with zeros.
